%% =======================================================
%%  Manuscript prepared for submission to
%%  Expert Systems with Applications (Elsevier)
%%
%%  STATUS:
%%   Numeric cells and factual placeholders have been filled from verified
%%   local experiment outputs. Ambiguous claims are framed conservatively.
%% =======================================================

\documentclass[preprint,12pt]{elsarticle}

\usepackage{amsmath}
\usepackage{amssymb}
\usepackage{graphicx}
\usepackage{booktabs}
\usepackage{multirow}
\usepackage{hyperref}
\usepackage{float}
\usepackage{microtype}
\usepackage{lineno}
\usepackage{xcolor}
\usepackage{diagbox}

\journal{Expert Systems with Applications}

\begin{document}

\begin{frontmatter}

%% =======================================================
%%  TITLE
%% =======================================================
\title{An intelligent vision-based decision-support system for
Parkinson's gait severity assessment using COM-centered
spatio-temporal modeling from a single RGB camera}

%% =======================================================
%%  AUTHORS
%% =======================================================
\author[cnu]{Bokyung Kim}
\author[cnu]{Hieyong Jeong\corref{cor1}}
\ead{h.jeong@jnu.ac.kr}
\author[cnu]{Md Ilias Bappi}
\author[cnu]{Kyungbaek Kim}

\cortext[cor1]{Corresponding author}

\address[cnu]{Chonnam National University, Gwangju, Republic of Korea}

%% =======================================================
%%  ABSTRACT
%% =======================================================
\begin{abstract}
Parkinson's disease (PD) is a progressive neurodegenerative disorder for
which gait impairment serves as a key clinical biomarker. In current
practice, gait severity is rated on the MDS-UPDRS Part~III scale through
expert visual judgement, which is inherently subjective, episodic, and
limited to clinical settings. An automated and deployable assessment tool
is therefore needed to support clinicians in objective decision making.

This paper proposes an intelligent decision-support system estimating
the MDS-UPDRS Part~III gait score directly from a single frontal-view RGB
video. The system integrates four key components: (i) a Centre-of-Mass
(COM)-centred coordinate normalisation scheme that removes global subject
translation; (ii) a multi-channel kinematic feature representation; (iii) a
Graph Convolutional Spatio-Temporal Encoder with learnable adjacency; and
(iv) bounded regression and selective prediction, which enforce the observable
clinical range and defer borderline cases to human review. We additionally
evaluate a median bone-length scale-normalised variant as a robustness
analysis, but the primary reported model is the COM-centred configuration.

We evaluate the system on a combined cohort of 131 PD patients from CNUH
($N=21$) and the multi-site CARE-PD subset ($N=110$, 6{,}066 sequences),
yielding 6{,}087 gait sequences in total. The main evaluation uses
subject-level GroupKFold cross-validation on the combined cohort; cross-site
zero-shot transfer in both directions is reported as an auxiliary external
generalisation analysis. Competing models are evaluated on the same
H36M-compatible 17-joint data and subject-level splits using model-specific
input adapters. Under the combined GroupKFold protocol, the proposed system
attains MAE\,=\,0.358 and the highest F1$_{0\text{--}2}$ among evaluated
models, while also providing bounded-output safety and an explicit clinical
abstention rule. Zero-shot transfer exposes a non-trivial domain gap between
CNUH and CARE-PD, whereas the bounded-regression head prevents catastrophic
extrapolation under severe low-data transfer. By additionally reporting
inference latency and a quantitative risk--coverage analysis of the abstention
rule, the proposed pipeline establishes a practical foundation for
privacy-preserving, edge-deployable PD gait monitoring.
\end{abstract}

%% =======================================================
%%  KEYWORDS
%% =======================================================
\begin{keyword}
Intelligent decision-support system \sep
Parkinson's disease \sep
Vision-based gait analysis \sep
Graph convolutional network \sep
Selective prediction \sep
MDS-UPDRS
\end{keyword}

\end{frontmatter}

\linenumbers

%% =======================================================
%%  HIGHLIGHTS  (paste into the separate Highlights file; each <=85 chars)
%% =======================================================
% * Vision-based decision-support system for MDS-UPDRS gait scoring from RGB video.
% * COM-centred modelling estimates MDS-UPDRS gait severity from RGB video.
% * Graph Convolutional Spatio-Temporal Encoder with bounded-output safety head.
% * Combined CNUH+CARE-PD evaluation with auxiliary zero-shot transfer tests.
% * Selective-prediction rule and edge latency reported for clinical deployment.

%% =======================================================
%%  1. INTRODUCTION
%% =======================================================
\section{Introduction}
\label{sec:intro}

\subsection{Background and clinical motivation}
Parkinson's disease (PD) is characterised by progressive motor symptoms
including bradykinesia, rigidity, and impaired gait \cite{poewe2017parkinson}.
Recent Global Burden of Disease estimates further indicate that the worldwide
PD burden has continued to increase over the past three decades
\cite{yang2026gbd}.
Gait severity is clinically rated using Part~III of the Movement Disorder
Society's Unified Parkinson's Disease Rating Scale (MDS-UPDRS)
\cite{goetz2008mds}. This rating relies on expert visual judgement,
exhibits non-trivial inter-rater variability \cite{richards1994inter}, and
is obtained only during scheduled visits, thereby missing the intraday
symptom fluctuations driven by medication cycles \cite{zhang2024wearable}.
These limitations motivate \emph{intelligent decision-support systems} that
can deliver objective, reproducible, and continuous quantification of gait
severity to assist clinicians.

\subsection{Why an RGB-only expert system?}
Marker-based optical motion-capture systems offer millimetre-level accuracy
but require dedicated laboratories \cite{dibiase2020gait}, while wearable
inertial sensors demand strict placement, charging, and skin-contact
adherence \cite{zhang2024wearable}. By contrast, an expert system that
estimates MDS-UPDRS gait scores directly from a single RGB camera meets
several practical deployment desiderata: zero hardware burden on the patient,
compatibility with existing clinic and consumer cameras, the potential for
privacy-preserving edge deployment when only skeletal keypoints are stored,
and scalability across sites \cite{sun2024digital}.

\subsection{Open problems and our contributions}
Despite progress in skeleton-based motion modelling
\cite{shi2019twoStream,zhu2023motionbert}, deploying PD gait assessment as a
reliable clinical system requires addressing three concrete gaps:

\begin{enumerate}
    \item \textbf{Robustness to environmental variation.} Predictions from a
    single RGB video are sensitive to camera-to-subject distance (apparent
    scale) and horizontal subject position; few algorithmic studies formulate
    and empirically quantify an explicit invariance mechanism against these
    spatial confounders \cite{gennarelli2025skelmamba}.
    \item \textbf{Evaluation bias in ordinal assessment.} Vision-based
    regression formulations frequently report Mean Absolute Error (MAE) alone,
    which can mask regression-to-the-mean behaviour on minority severe cases;
    the classification benchmarks of this field instead report Macro-F1
    \cite{adeli2025carepd}.
    \item \textbf{Clinical decision support versus black-box prediction.}
    A fielded expert system must provide deployment transparency, operational
    rules such as bounded clinical ranges, abstention mechanisms for uncertain
    cases, and measured computational cost \cite{sun2024digital}.
\end{enumerate}

To address these gaps, we propose an intelligent decision-support system with
the following contributions:
\begin{itemize}
    \item We formalise a \textbf{COM-centred spatial normalisation scheme}
    that removes global subject translation and evaluate an optional
    median bone-length scaling step as a robustness variant for camera-distance
    perturbations (Section~\ref{sec:robust_exp}).
    \item We process multi-channel kinematic features with a \textbf{Graph
    Convolutional Spatio-Temporal Encoder} featuring learnable adjacency
    \cite{shi2019twoStream}, discovering non-anatomical inter-joint
    dependencies indicative of pathological gait compensations.
    \item We evaluate rigorously with both \textbf{MAE and Macro-F1}
    (reported as F1$_{0\text{--}3}$ and F1$_{0\text{--}2}$ to match the CARE-PD
    benchmark), transparently handling class imbalance, and we run every
    competing model under the same \textbf{H36M-compatible 17-joint
    representation and subject-level splits}, using only model-specific input
    adapters required by each architecture.
    \item We explicitly evaluate \textbf{zero-shot cross-site generalisation}
    in both directions (CARE-PD\,$\leftrightarrow$\,CNUH), revealing the
    expected domain gap while showing that the \textbf{bounded-regression}
    safety head prevents the catastrophic extrapolation seen in unbounded
    baselines under severe low-data transfer.
    \item We formalise a \textbf{selective-prediction (abstention) rule},
    quantify it with a risk--coverage analysis against a random-abstention
    baseline, and report practical \textbf{edge-inference latency}
    (Section~\ref{sec:discussion}).
\end{itemize}

%% =======================================================
%%  2. RELATED WORK
%% =======================================================
\section{Related Work}
\label{sec:related}

\subsection{Vision-based gait analysis for PD}
The maturity of markerless pose estimation \cite{cao2017openpose,
bazarevsky2020blazepose} has enabled non-contact PD gait analysis from RGB
video. Lu et al. \cite{lu2020vision,lu2021quantifying} pioneered vision-based
MDS-UPDRS gait scoring from 3D pose sequences. The CARE-PD benchmark
\cite{adeli2025carepd} consolidates 363 participants across nine cohorts from
eight clinical sites, distributed as harmonised SMPL meshes, and defines
standardised generalisation protocols (within-dataset, cross-dataset,
leave-one-dataset-out, and in-domain adaptation), enabling large-scale
validation of motion encoders \cite{gennarelli2025skelmamba,
tian2024crossspatiotemporal,adeli2024benchmarking}. We utilise a curated
110-participant, fully annotated subset of CARE-PD as our primary training
domain. A recurring limitation across this literature is the unquantified
sensitivity of predictions to camera scale and subject distance.

\subsection{Skeleton-based deep models and expert systems}
Skeleton modelling has evolved from fixed-graph convolutions
(ST-GCN \cite{yan2018spatial}) to adaptive graphs
(2s-AGCN \cite{shi2019twoStream}) and Transformer-based encoders
(MotionBERT \cite{zhu2023motionbert},
MotionAGFormer \cite{mehraban2024motionagformer}). From an expert-systems
perspective, however, raw accuracy is insufficient: a fielded system must also
exhibit domain-specific robustness, interpretable failure modes, an explicit
uncertainty/abstention policy, and edge-level efficiency. Our work bridges
this gap by embedding a biomechanically motivated normalisation and
deployment-centric decision rules into a modern spatio-temporal architecture.

%% =======================================================
%%  3. METHODOLOGY
%% =======================================================
\section{Proposed System}
\label{sec:method}

The proposed system takes a single frontal-view RGB gait video and outputs a
continuous MDS-UPDRS Part~III (Item~3.10) gait estimate. It comprises four
stages: scale-robust normalisation, multi-channel feature construction,
spatio-temporal encoding, and a bounded-regression head.

\subsection{Stage 1: Pose extraction and COM-centred normalisation}
\label{sec:stage1}

\paragraph{Pose estimation.}
We apply the MediaPipe Pose estimator \cite{bazarevsky2020blazepose} to each
frame, obtaining per-joint image-plane coordinates and a model-predicted
relative depth. We treat the result as a 2.5D coordinate
$\mathbf{p}^{\text{raw}}_{t,j}\in\mathbb{R}^{3}$, used as structured input
rather than as a calibrated 3D measurement, with $T=390$ frames and $J=17$
H36M-compatible joints. For CARE-PD, we used the released SMPL-derived
H36M-compatible skeleton representation. For CNUH, MediaPipe 33-landmark
sequences were mapped to the same H36M-compatible 17-joint layout, yielding a
common input representation across cohorts.

\paragraph{COM-centred normalisation.}
To remove the subject's global image-plane position, we approximate the body
centre (COM proxy) as the midpoint of the left and right hips, which is
biomechanically co-located with the trunk's centre of mass during quiet
bipedal locomotion \cite{winter2009biomechanics}:
\begin{equation}
\mathbf{C}_t = \tfrac{1}{2}\!\left(
\mathbf{p}^{\text{raw}}_{t,\text{L-hip}}+\mathbf{p}^{\text{raw}}_{t,\text{R-hip}}\right),
\qquad
\mathbf{p}^{\text{COM}}_{t,j}=\mathbf{p}^{\text{raw}}_{t,j}-\mathbf{C}_t .
\label{eq:trans}
\end{equation}
The primary proposed model uses $\mathbf{p}^{\text{COM}}_{t,j}$ as its
coordinate stream; this is the configuration corresponding to the main
MAE\,=\,0.358 result.

\paragraph{Scale-normalised robustness variant.}
For the robustness experiment only, we additionally divide the COM-centred
coordinates by the median torso bone length $L_{\text{ref}}$ computed over the
sequence:
\begin{equation}
\mathbf{p}^{\text{scale}}_{t,j}=\frac{\mathbf{p}^{\text{COM}}_{t,j}}{L_{\text{ref}}}.
\label{eq:norm}
\end{equation}
Under a global multiplicative model of distance change, this variant is
invariant to horizontal translation and coordinate scaling by construction;
Section~\ref{sec:robust_exp} verifies this empirically and discusses its limits
under true perspective distortion.

\subsection{Stage 2: Multi-channel kinematic feature construction}
\label{sec:stage2}
From the selected coordinate tensor (COM-centred for the main model and
COM+scale-normalised for the robustness variant), we construct a per-(frame,
joint) feature vector by concatenating clinically motivated channels---joint positions
($f_{\text{pos}}$), frame-difference velocities ($f_{\text{vel}}$,
sensitive to bradykinesia), sequence-level motion amplitude and variability
($f_{\text{amp}},f_{\text{var}}$), and joint angles
($f_{\text{ang}}$, encoding range of motion):
\begin{equation}
\mathbf{x}_{t,j}=\mathrm{Concat}\!\left[
f_{\text{pos}},f_{\text{vel}},f_{\text{amp}},f_{\text{var}},f_{\text{ang}}\right]\in\mathbb{R}^{C}.
\end{equation}

\subsection{Stage 3: Graph Convolutional Spatio-Temporal Encoder}
\label{sec:stage3}
Unlike ST-GCN \cite{yan2018spatial}, which relies on a fixed anatomical graph,
we use a \emph{learnable} adjacency matrix \cite{shi2019twoStream} to discover
task-relevant, non-anatomical inter-joint dependencies characteristic of
pathological gait compensation. Two GraphConv layers (with LayerNorm and ReLU)
expand the per-joint dimension ($C\!\to\!64\!\to\!128$). A multi-head
self-attention module over the $J=17$ joint tokens captures global spatial
relations, and a Temporal Transformer encoder along the time axis models
frame-to-frame dynamics including stride periodicity. The sequence is collapsed
by temporal average pooling to $\mathbf{h}\in\mathbb{R}^{J\times128}$.

\subsection{Stage 4: Bounded regression head}
\label{sec:stage4}
To enforce clinical validity, the pooled representation is mapped by an MLP to a
sigmoid-scaled bounded output:
\begin{equation}
\hat{y}=3\cdot\sigma\!\left(\mathrm{MLP}(\mathbf{h})\right)\in(0,3).
\label{eq:bounded}
\end{equation}
Because score~4 denotes inability to walk independently---which precludes
frontal-view ambulatory video---the range $[0,3]$ covers all feasible cases in
both cohorts. The network is trained by minimising the MAE between $\hat{y}$ and
the ground-truth score $y$.

\subsection{Implementation and reproducibility}
\label{sec:repro}
Table~\ref{tab:hparams} reports all hyperparameters. Source code and
configuration files are released at
\url{https://github.com/bokyung08/parkinson_COM}.

\begin{table}[htpb]
\centering
\caption{Implementation and training hyperparameters of the proposed system.}
\label{tab:hparams}
\resizebox{\linewidth}{!}{%
\begin{tabular}{ll}
\toprule
\textbf{Component} & \textbf{Setting} \\
\midrule
Pose estimator               & MediaPipe Pose (BlazePose Heavy) \\
Joints / Sequence length     & $J=17$ (H36M-compatible) / $T=390$ (interpolation) \\
Normalisation                & COM translation (main); median bone scaling only for robustness variant \\
GraphConv hidden dim         & $C \to 64 \to 128$ (learnable adjacency) \\
Joint-attention heads        & 4 \\
Temporal Transformer         & 1 layer, 4 heads; no explicit positional encoding \\
MLP hidden dim / Dropout     & 256 / 0.1 \\
Output activation            & $3\cdot\sigma(\cdot)$ (bounded regression) \\
Total parameters             & 158{,}594 \\
\midrule
Optimiser / LR / Batch       & Adam / $1\times10^{-4}$ / 16 \\
Epochs / Loss                & 100 / MAE \\
Framework / Hardware         & PyTorch 2.0 / NVIDIA RTX 3080 \\
\bottomrule
\end{tabular}%
}
\end{table}

%% =======================================================
%%  4. EXPERIMENTAL SETUP
%% =======================================================
\section{Experimental Setup}
\label{sec:setup}

\subsection{Datasets and cohort roles}
\label{sec:datasets}
\paragraph{CARE-PD subset.}
We use a fully annotated subset of the CARE-PD benchmark
\cite{adeli2025carepd} comprising 110 participants and 6{,}066 sequences drawn
from four cohorts (3DGait, BMCLab, PD-GaM, T-SDU-PD).
Although CARE-PD contains nine source cohorts in total, the current H36M17
MDS-UPDRS gait-score experiments use the four cohorts available in the local
processed manifest with compatible item-10 labels (Table~\ref{tab:carepd_split}).
All GroupKFold and LODO results reported in this manuscript are constructed
from these four CARE-PD cohorts plus the CNUH cohort where explicitly stated.

\begin{table}[htpb]
\centering
\caption{CARE-PD cohorts included in the processed H36M17 item-10 manifest.
All four cohorts are included in GroupKFold and CARE-PD LODO experiments.}
\label{tab:carepd_split}
\resizebox{\linewidth}{!}{%
\begin{tabular}{lrrrrrr}
\toprule
\textbf{Cohort} & \textbf{Seq.} & \textbf{Subj.} &
\textbf{Score 0} & \textbf{Score 1} & \textbf{Score 2} & \textbf{Score 3} \\
\midrule
3DGait   &   90 & 43 &  24 &  42 &  14 & 10 \\
BMCLab   & 3895 & 23 & 1705 & 1380 & 810 &  0 \\
PD-GaM   & 1700 & 30 &  783 &  635 & 248 & 34 \\
T-SDU-PD &  381 & 14 &   96 &  118 & 167 &  0 \\
\midrule
\textbf{Total} & \textbf{6066} & \textbf{110} &
\textbf{2608} & \textbf{2175} & \textbf{1239} & \textbf{44} \\
\bottomrule
\end{tabular}%
}
\end{table}

\paragraph{CNUH clinical cohort.}
Frontal-view gait videos were prospectively collected from $N=21$ PD patients
at Chonnam National University Hospital (IRB: CNUH-2025-203; written informed
consent obtained).

\paragraph{Main combined cohort and transfer cohorts.}
The main reported GroupKFold results are computed on the combined
CNUH+CARE-PD cohort, containing 131 subjects and 6{,}087 gait sequences. To
separately assess cross-site transferability, we also train on one cohort and
evaluate on the other without fine-tuning or domain adaptation.

\subsection{Evaluation protocol and metrics}
\label{sec:metrics}
We adopt subject-level GroupKFold 5-fold cross-validation on the combined
CNUH+CARE-PD cohort, ensuring that all sequences of a subject lie in a single
fold. We report MAE and
RMSE in score units, and Macro-F1 obtained by rounding the continuous output to
the nearest integer (clipped to $[0,3]$). Following the CARE-PD benchmark, we
report two Macro-F1 variants: F1$_{0\text{--}3}$ over all classes and
F1$_{0\text{--}2}$ excluding the rare score-3 class (the metric only; all
score-3 samples remain in training). We note that rounding a regressor is not
identical to a native ordinal classifier, so Macro-F1 comparisons with
classification benchmarks are indicative rather than exact. For significance
testing we use the Wilcoxon signed-rank test. To avoid inflating significance by
treating multiple sequences from the same subject as independent samples,
absolute errors were also aggregated at the subject level before testing,
yielding $p=0.107$ for Ours versus Lu et al. and
$p=9.09\times10^{-4}$ for Ours versus ST-GCN.

\subsection{Fair-comparison protocol}
\label{sec:fairness}
All competing models---classical regressors, ST-GCN \cite{yan2018spatial},
Lu et al.\ \cite{lu2020vision}, MotionBERT \cite{zhu2023motionbert}, and
MotionAGFormer \cite{mehraban2024motionagformer}---are trained and evaluated on
the same H36M-compatible 17-joint representation and the same subject-level
GroupKFold split. The proposed model uses the full Configuration~D feature
stream, whereas each external architecture uses the input adapter required by
its original design. Base sequences are resampled to $T=390$, and
encoder-specific temporal windows are used where required by pretrained motion
encoders. Publicly released checkpoints were used for initialisation where
available. Consequently, the comparison in Section~\ref{sec:comparison}
reflects performance under a common data protocol rather than a claim of
identical internal input tensors across all architectures.

\paragraph{Lu et al. reimplementation.}
Lu et al. originally estimated MDS-UPDRS scores from 3D pose sequences
extracted from video. Because the present experiments use a unified
H36M-compatible 17-joint representation across CNUH and CARE-PD, we preserve
the released DD-Net/OF-DDNet architecture and adapt only the input interface to
our 17-joint coordinate stream and derived motion features. The model is trained
on the same subject-level splits and target labels as all other baselines, with
no additional dataset-specific fine-tuning. We therefore interpret this result
as an official-architecture Lu et al. baseline under the present 17-joint
2.5D/SMPL-derived input setting, not as an exact reproduction of the original
MICCAI 2020 3D-pose experimental condition.

%% =======================================================
%%  5. EXPERIMENTAL RESULTS
%% =======================================================
\section{Experimental Results}
\label{sec:results}

\subsection{Quantitative robustness to spatial perturbations}
\label{sec:robust_exp}
To validate Eq.~\eqref{eq:norm}, we apply synthetic scale (distance) and
horizontal-translation perturbations to the input coordinates at inference time
(Table~\ref{tab:robustness}).

\begin{table}[htpb]
\centering
\caption{Robustness to affine scale (distance) and translation perturbations.
$\Delta$MAE is the relative change from each model's own unperturbed baseline.}
\label{tab:robustness}
\resizebox{\linewidth}{!}{%
\begin{tabular}{lcc}
\toprule
\textbf{Perturbation} & \textbf{w/o Bone-norm ($\Delta$MAE)} &
\textbf{Bone-norm variant ($\Delta$MAE)} \\
\midrule
Original (baseline MAE)        & \textbf{0.358} & 0.366 \\
Scale $s=0.70$ (far)           & $+35.9\%$ & \textbf{$+0.0\%$} \\
Scale $s=1.30$ (near)          & $+46.5\%$ & \textbf{$+0.0\%$} \\
Shift $\Delta x=\pm0.20$       & $+0.0\%$  & \textbf{$+0.0\%$} \\
\bottomrule
\end{tabular}%
}
\end{table}

Without bone-length scaling, distance-induced scale perturbations degrade MAE by
up to $46.5\%$. Adding the scaling step eliminates this shift entirely
($\Delta$MAE $=0\%$). This invariance holds \emph{by construction} under a global
multiplicative model of distance change, and we do not claim robustness to true
perspective distortion arising from extreme camera angles, which we treat as a
limitation (Section~\ref{sec:failure}). The marginal accuracy cost of adding the
step (0.358\,$\to$\,0.366 MAE) is justified when explicit scale invariance is
required; we therefore report the scale-normalised configuration as a robustness
variant while using the 0.358-MAE COM-centred configuration as the main
proposed model.

\subsection{Ablation studies}
\label{sec:ablation}
Table~\ref{tab:ablation_feat} reports the input-feature ablation and
Table~\ref{tab:ablation_arch} the architecture ablation, both on the
COM-centred combined CNUH+CARE-PD cohort under the GroupKFold protocol.

\begin{table}[htpb]
\centering
\caption{Input-feature ablation on the combined CNUH+CARE-PD cohort
(GroupKFold 5-fold, $N=6{,}087$).
Lower MAE/RMSE and higher F1 are better.}
\label{tab:ablation_feat}
\resizebox{\linewidth}{!}{%
\begin{tabular}{clccc}
\toprule
\textbf{Config.} & \textbf{Feature set} & \textbf{MAE} & \textbf{RMSE} &
\textbf{F1$_{0\text{--}3}$} \\
\midrule
A & COM coordinates only            & 0.374 & 0.577 & 0.617 \\
B & + velocity                      & 0.432 & 0.629 & 0.508 \\
C & + amplitude / variability       & 0.376 & 0.549 & 0.546 \\
D & \textbf{+ joint angle (Proposed)} & \textbf{0.358} & 0.564 & \textbf{0.661} \\
\bottomrule
\end{tabular}%
}
\end{table}

\begin{table}[htpb]
\centering
\caption{Architecture ablation on the combined CNUH+CARE-PD cohort
(Configuration~D features, GroupKFold 5-fold).}
\label{tab:ablation_arch}
\resizebox{\linewidth}{!}{%
\begin{tabular}{lccc}
\toprule
\textbf{Model} & \textbf{MAE} & \textbf{RMSE} & \textbf{F1$_{0\text{--}3}$} \\
\midrule
MLP only (mean pooling)            & 0.554 & 0.653 & 0.381 \\
GraphConv + MLP                    & 0.450 & 0.580 & 0.510 \\
GraphConv + Joint Attention + MLP  & 0.414 & 0.564 & 0.507 \\
\textbf{Full model (Ours)}         & \textbf{0.358} & 0.564 & \textbf{0.661} \\
\bottomrule
\end{tabular}%
}
\end{table}

\subsection{Comparison with state-of-the-art motion encoders}
\label{sec:comparison}
Table~\ref{tab:comparison} compares the proposed system against classical,
graph-based, and Transformer-based baselines, all under the fair-comparison
protocol of Section~\ref{sec:fairness}.

\begin{table}[htpb]
\centering
\caption{Comparison on the combined CNUH+CARE-PD cohort (17-joint,
subject-level GroupKFold 5-fold, $N=6{,}087$). All models use the same
subject-level split and H36M-compatible skeleton representation with
model-specific input adapters.}
\label{tab:comparison}
\resizebox{\linewidth}{!}{%
\begin{tabular}{lccccc}
\toprule
\textbf{Model} & \textbf{\#Params} & \textbf{MAE} & \textbf{RMSE} &
\textbf{F1$_{0\text{--}3}$} & \textbf{F1$_{0\text{--}2}$} \\
\midrule
SVR (best classical)              & ---     & 0.492 & 0.639 & 0.485 & 0.597 \\
Lu et al.\ \cite{lu2020vision}    & 147{,}908 & 0.404 & 0.543 & \textbf{0.666} & 0.673 \\
ST-GCN \cite{yan2018spatial}      & 252{,}097 & 0.443 & 0.623 & 0.476 & 0.628 \\
MotionBERT \cite{zhu2023motionbert}        & 10.8M & 0.442 & 0.625 & 0.457 & 0.613 \\
MotionAGFormer \cite{mehraban2024motionagformer} & 2.3M & 0.405 & 0.638 & 0.561 & 0.636 \\
\midrule
\textbf{Proposed System (Ours)}   & 158{,}594 & \textbf{0.358} & 0.564 & 0.661 & \textbf{0.691} \\
\bottomrule
\end{tabular}%
}
\end{table}

With only 158{,}594 parameters, the proposed system performs competitively with
substantially larger Transformer encoders in both MAE and Macro-F1 while adding
bounded-output safety and a scale-normalised robustness variant. Because all
baselines are evaluated under the same subject-level split and 17-joint
skeleton representation, the comparison supports an end-to-end model-level
assessment under a common data protocol.
The proposed system attains the lowest MAE (0.358) and the highest
F1$_{0\text{--}2}$ (0.691). Lu et al. attains a marginally higher
F1$_{0\text{--}3}$ (0.666 vs. 0.661) and lower RMSE, while MotionAGFormer
attains lower MedAE in the supplementary analysis. We therefore frame the final
claim around the primary MAE metric, score-0--2 Macro-F1, bounded-output safety,
selective prediction, and computational efficiency rather than claiming
dominance on every metric.

\subsection{Cross-site zero-shot transfer and bounded safety}
\label{sec:crossds}

\paragraph{Auxiliary zero-shot transfer (CARE-PD\,$\to$\,CNUH).}
To assess cross-site transfer, the model is trained on the CARE-PD subset
($N=110$) and evaluated, without any fine-tuning, on the independent CNUH cohort
($N=21$) (Table~\ref{tab:zero_primary}).

\begin{table}[htpb]
\centering
\caption{Auxiliary zero-shot transfer (CARE-PD\,$\to$\,CNUH). No fine-tuning or
domain adaptation was applied.}
\label{tab:zero_primary}
\resizebox{\linewidth}{!}{%
\begin{tabular}{lccc}
\toprule
\textbf{Model} & \textbf{MAE} & \textbf{RMSE} & \textbf{F1$_{0\text{--}3}$} \\
\midrule
ST-GCN \cite{yan2018spatial}      & 1.203 & 1.385 & 0.100 \\
Lu et al.\ \cite{lu2020vision}    & \textbf{0.865} & \textbf{1.027} & 0.155 \\
\textbf{Proposed (Ours)}          & 0.910 & 1.034 & \textbf{0.178} \\
\bottomrule
\end{tabular}%
}
\end{table}

\paragraph{Extreme low-data stress test (CNUH\,$\to$\,CARE-PD).}
To isolate the effect of the bounded-regression head, we reverse the direction:
training on only 21 CNUH subjects and extrapolating to 6{,}066 CARE-PD sequences
(Table~\ref{tab:zero_stress}). This is a deliberate stress test of numerical
stability under severe data scarcity, not a deployment scenario.

\begin{table}[htpb]
\centering
\caption{Extreme low-data stress test (CNUH\,$\to$\,CARE-PD), isolating the
bounded-output safety mechanism.}
\label{tab:zero_stress}
\resizebox{\linewidth}{!}{%
\begin{tabular}{lcc}
\toprule
\textbf{Model variant} & \textbf{MAE} & \textbf{RMSE} \\
\midrule
ST-GCN (unbounded)            & 8.346 & 9.737 \\
ST-GCN (+ bounded head)       & 2.072 & 2.201 \\
\textbf{Proposed (bounded)}   & \textbf{0.747} & \textbf{0.882} \\
\bottomrule
\end{tabular}%
}
\end{table}

The unbounded ST-GCN suffers catastrophic failure (MAE 8.346) from extreme
extrapolation on domain-shifted data. Adding the same bounded output range
reduces this failure substantially (MAE 2.072), confirming the value of bounded
output for numerical safety. However, the bounded ST-GCN remains far worse than
the proposed model (MAE 0.747), indicating that output bounding alone is
insufficient and that the proposed graph-attention-temporal architecture is also
needed for low transfer error.

\subsection{Clinical performance and transparency}
\label{sec:perclass}
Table~\ref{tab:perclass} stratifies error by ground-truth score, and
Table~\ref{tab:lodo} reports the CARE-PD leave-one-dataset-out (LODO)
protocol, a cohort-level generalisation test.

\begin{table}[htpb]
\centering
\caption{Per-class error and per-class F1 on the combined CNUH+CARE-PD cohort
(GroupKFold 5-fold, $N=6{,}087$).}
\label{tab:perclass}
\resizebox{\linewidth}{!}{%
\begin{tabular}{crccc}
\toprule
\textbf{Score} & \textbf{$N$ (seq.)} & \textbf{MAE} & \textbf{RMSE} &
\textbf{Class F1} \\
\midrule
0 & 2{,}615 & 0.225 & 0.395 & 0.793 \\
1 & 2{,}183 & 0.342 & 0.482 & 0.662 \\
2 & 1{,}244 & 0.649 & 0.889 & 0.612 \\
3 &    45  & 0.738 & 0.930 & 0.576 \\
\bottomrule
\end{tabular}%
}
\end{table}
The row-normalised confusion matrix and calibration-style curve are provided as
supplementary transparency figures.

\begin{table}[htpb]
\centering
\caption{CARE-PD leave-one-dataset-out (LODO) results. Each row holds out one
source cohort; Overall aggregates all held-out sequences.}
\label{tab:lodo}
\resizebox{\linewidth}{!}{%
\begin{tabular}{lccc}
\toprule
\textbf{Held-out cohort} & \textbf{MAE} & \textbf{RMSE} &
\textbf{F1$_{0\text{--}3}$} \\
\midrule
3DGait    & 0.775 & 0.947 & 0.233 \\
BMCLab    & 0.663 & 0.844 & 0.340 \\
PD-GaM    & 0.495 & 0.724 & 0.414 \\
T-SDU-PD  & 0.692 & 0.836 & 0.305 \\
\midrule
\textbf{Overall} & \textbf{0.620} & \textbf{0.813} & \textbf{0.358} \\
\bottomrule
\end{tabular}%
}
\end{table}

%% =======================================================
%%  6. DISCUSSION: SYSTEM DEPLOYMENT
%% =======================================================
\section{Discussion: A Clinical Decision-Support System}
\label{sec:discussion}

To qualify as an expert system rather than a bare predictor, the model must
manage its own uncertainty and report concrete deployment cost. We address both.

\subsection{Selective prediction: a risk--coverage analysis}
\label{sec:selective}
A deployable clinical system should defer unreliable predictions to a human
expert. We equip the system with a \emph{post-hoc abstention rule} requiring no
retraining. Because the rounded decision boundaries of the ordinal scale lie at
the half-integers $\{0.5,1.5,2.5\}$, the distance of a continuous prediction
$\hat{y}$ to its nearest boundary is a natural confidence proxy:
\begin{equation}
u(\hat{y}) = \min_{b\in\{0.5,1.5,2.5\}} \lvert \hat{y}-b\rvert .
\label{eq:uncertainty}
\end{equation}
A small $u$ indicates a boundary-proximal prediction most likely to be misrounded;
a large $u$ indicates a prediction comfortably inside a class interval. We
\emph{abstain} (flag for clinician review) whenever $u(\hat{y})<\tau$ and
auto-score the remainder. Sweeping $\tau$ traces a \emph{risk--coverage curve},
where \emph{coverage} is the fraction of cases auto-scored and \emph{selective
risk} is the MAE on the retained subset.

Discarding hard cases trivially lowers retained error, so to demonstrate that
the gain reflects an \emph{informative} uncertainty signal we compare against a
\emph{random-abstention} baseline that flags the same fraction of cases
uniformly at random; an informative rule must lie strictly below this baseline
at equal coverage. We summarise the full curve by the Area Under the
Risk--Coverage curve (AURC; lower is better). We additionally report a fixed,
clinically interpretable operating band $\hat{y}\in[1.3,1.7]$, which targets the
contentious $1\!\leftrightarrow\!2$ boundary where inter-rater disagreement is
highest \cite{richards1994inter}.

\begin{table}[htpb]
\centering
\caption{Selective prediction (risk--coverage) on the combined CNUH+CARE-PD
cohort. At each
target coverage, the threshold $\tau$ flags the most boundary-proximal cases via
Eq.~\eqref{eq:uncertainty}. We report MAE and F1$_{0\text{--}3}$ on the retained
(auto-scored) subset, and the MAE of a random-abstention baseline at matched
coverage. The final row is the fixed clinical band $[1.3,1.7]$.}
\label{tab:selective}
\resizebox{\linewidth}{!}{%
\begin{tabular}{lccccc}
\toprule
\textbf{Coverage} & \textbf{$\tau$} & \textbf{Flagged (\%)} &
\textbf{Retained MAE} & \textbf{Retained F1$_{0\text{--}3}$} &
\textbf{Random-abst.\ MAE} \\
\midrule
100\% (no abstention) & ---          & 0.0           & 0.358        & 0.661 & 0.358 \\
90\%                  & 0.11  & 10.0          & 0.325 & 0.687 & 0.358 \\
80\%                  & 0.22  & 20.0          & 0.288 & 0.692 & 0.357 \\
70\%                  & 0.30  & 30.0          & 0.260 & 0.701 & 0.358 \\
\midrule
Clinical band $[1.3,1.7]$ & ---      & 6.0    & 0.345 & 0.675 & 0.358 \\
\midrule
\multicolumn{6}{l}{AURC (Ours): 0.278 \quad vs.\ random abstention:
0.358 \quad (lower is better)} \\
\bottomrule
\end{tabular}%
}
\end{table}

As Table~\ref{tab:selective} shows, retained MAE decreases from 0.358 at full
coverage to 0.260 at 70\% coverage (a 27.2\% relative reduction),
while the random-abstention baseline at the same coverage remains essentially
flat (0.358). The gap confirms that the boundary-distance score of
Eq.~\eqref{eq:uncertainty} is an informative uncertainty proxy rather than a
trivial discarding of hard samples. The fixed clinical band $[1.3,1.7]$ flags
only 6.0\% of cases yet reduces retained MAE to 0.345, giving a simple,
auditable human-in-the-loop operating point; the AURC of 0.278 versus
0.358 for random abstention quantifies the benefit over the entire curve.

\subsection{Edge deployability}
\label{sec:deploy}
With 158{,}594 parameters, the model requires 0.242\,ms per sample on an RTX
3080 under the CUDA batch-32 latency benchmark and 5.793\,ms per sample on the
measured CPU benchmark, enabling local inference without cloud infrastructure.
Once skeletal keypoints are extracted, the raw RGB stream is
discarded, so only low-dimensional joint coordinates are stored or transmitted,
substantially reducing re-identification risk relative to pixel-based pipelines.

\subsection{Failure modes and limitations}
\label{sec:failure}
\begin{itemize}
    \item \emph{Class imbalance.} Score-3 sequences are severely
    underrepresented ($\sim$0.7\%), limiting discriminative capacity for the most
    severe cases; class-weighted losses or targeted augmentation are natural
    extensions.
    \item \emph{Perspective distortion.} COM normalisation neutralises global
    translation, and the optional bone-length scaling variant neutralises global
    2D scale \emph{by construction}; neither step models full 3D perspective
    distortion from extreme camera angles. Robustness under real
    multi-distance, multi-angle capture remains for dedicated evaluation.
    \item \emph{Pose-extraction failure.} Under heavy occlusion MediaPipe may lose
    tracking, drifting predictions toward the cohort mean; confidence-based frame
    filtering is recommended for uncontrolled environments.
    \item \emph{Statistical testing.} Sequence-level significance tests inflate
    the effective sample size because sequences cluster within subjects;
    confirmatory analysis aggregates errors at the subject level
    (Section~\ref{sec:metrics}).
\end{itemize}

%% =======================================================
%%  7. CONCLUSION
%% =======================================================
\section{Conclusion}
\label{sec:conclusion}
We presented an intelligent vision-based decision-support system for Parkinson's
gait severity assessment from a single RGB camera. Beyond a bare motion encoder,
the system addresses practical deployment bottlenecks: it provides closed-form
translation handling via COM-centred normalisation, an optional
scale-normalised robustness variant under a global-scaling approximation,
clinical output bounds, and an explicit abstention rule quantified by a
risk--coverage analysis. Validated on the combined CNUH+CARE-PD cohort under
subject-level GroupKFold cross-validation, the system achieves the best MAE
(0.358) and the best F1$_{0\text{--}2}$ among evaluated encoders at a fraction
of their parameter count. Auxiliary zero-shot transfer experiments reveal a
substantial cross-site domain gap, while bounded regression improves numerical
safety under severe low-data transfer. By reporting edge-inference latency,
Macro-F1 under class imbalance, and transparent failure modes, the pipeline
forms a practical, privacy-preserving building block for clinical decision
support in Parkinson's disease.

%% =======================================================
%%  CRediT
%% =======================================================
\section*{CRediT authorship contribution statement}
\textbf{Bokyung Kim}: Conceptualisation, Methodology, Software, Formal analysis,
Data curation, Writing -- original draft, Visualisation.
\textbf{Hieyong Jeong}: Conceptualisation, Supervision, Writing -- review and
editing, Funding acquisition.
\textbf{Md Ilias Bappi}: Investigation, Writing -- review and editing.
\textbf{Kyungbaek Kim}: Resources, Writing -- review and editing.

\section*{Declaration of competing interest}
The authors declare no known competing financial interests or personal
relationships that could have influenced the work reported in this paper.

\section*{Data availability}
Source code is available at
\url{https://github.com/bokyung08/parkinson_COM}. The CNUH clinical data cannot
be redistributed due to patient-privacy regulations; requests may be directed to
the corresponding author. CARE-PD is available from \cite{adeli2025carepd}.

\section*{Ethical approval}
This study was approved by the IRB of Chonnam National University Hospital
(CNUH-2025-203). Written informed consent was obtained from all participants, in
accordance with the 1964 Helsinki Declaration and its later amendments.

\section*{Acknowledgements}
This work was supported by the National Research Foundation of Korea (NRF) grant
funded by the Korea government (MSIT) (No.~RS-2026-25469002), and by the
IITP-ITRC grant funded by the Korea government (MSIT)
(IITP-2026-RS-2024-00437718).

%% =======================================================
%%  REFERENCES
%% =======================================================
\bibliographystyle{elsarticle-num}
\begin{thebibliography}{99}
\bibitem{poewe2017parkinson} Poewe W, Seppi K, Tanner CM, Halliday GM, Brundin P, Volkmann J, Schrag AE, Lang AE. Parkinson disease. \textit{Nat Rev Dis Primers} 2017;3:17013. doi:10.1038/nrdp.2017.13.
\bibitem{yang2026gbd} Yang R, Sun M, Chen W, Feng H, Chen B, Liu Y, He Q, Wang L, Zou C, Luo X, Li Z, Fu A, Qiao F, Tang H, Yang J, Ren H. Global, regional and national burden of Parkinson's disease, 1990--2021: Update from the GBD 2021 study. \textit{J Neurol Sci} 2026;480:125703. doi:10.1016/j.jns.2025.125703.
\bibitem{richards1994inter} Richards M, et al. Inter-rater reliability of the Unified Parkinson's Disease Rating Scale. \textit{Mov Disord} 1994;9(1):89--91.
\bibitem{goetz2008mds} Goetz CG, et al. MDS-UPDRS: scale presentation and clinimetric testing results. \textit{Mov Disord} 2008;23(15):2129--2170.
\bibitem{dibiase2020gait} di Biase L, et al. Gait analysis in Parkinson's disease: most accurate markers for diagnosis and monitoring. \textit{Sensors} 2020;20(12):3529.
\bibitem{zhang2024wearable} Zhang W, et al. Wearable sensor-based quantitative gait analysis in Parkinson's disease. \textit{npj Digit Med} 2024;7:169.
\bibitem{sun2024digital} Sun Y, et al. Digital biomarkers for precision diagnosis and monitoring in Parkinson's disease. \textit{npj Digit Med} 2024;7:218.
\bibitem{cao2017openpose} Cao Z, et al. OpenPose: realtime multi-person 2D pose estimation using part affinity fields. \textit{IEEE TPAMI} 2019;43(1):172--186.
\bibitem{bazarevsky2020blazepose} Bazarevsky V, et al. BlazePose: on-device real-time body pose tracking. \textit{arXiv} 2020;arXiv:2006.10204.
\bibitem{lu2020vision} Lu M, et al. Vision-based estimation of MDS-UPDRS gait scores for assessing Parkinson's disease motor severity. In: \textit{MICCAI}; 2020. p.\ 637--647.
\bibitem{lu2021quantifying} Lu M, et al. Quantifying Parkinson's disease motor severity under uncertainty using MDS-UPDRS videos. \textit{Med Image Anal} 2021;73:102179.
\bibitem{gennarelli2025skelmamba} Gennarelli I, et al. SkelMamba: a markerless motion capture approach for automated diagnosis of neurological disorders based on RGB videos. \textit{Gait Posture} 2025;122:109944.
\bibitem{tian2024crossspatiotemporal} Tian H, et al. Cross-spatiotemporal graph convolution networks for skeleton-based Parkinsonian gait MDS-UPDRS score estimation. \textit{IEEE TNSRE} 2024;32:412--421.
\bibitem{adeli2025carepd} Adeli V, et al. CARE-PD: a multi-site anonymized clinical dataset for Parkinson's disease gait assessment. In: \textit{NeurIPS Datasets and Benchmarks}; 2025. arXiv:2510.04312. doi:10.48550/arXiv.2510.04312.
\bibitem{adeli2024benchmarking} Adeli V, et al. Benchmarking skeleton-based motion encoder models for clinical applications: estimating Parkinson's disease severity in walking sequences. In: \textit{IEEE FG}; 2024.
\bibitem{yan2018spatial} Yan S, Xiong Y, Lin D. Spatial temporal graph convolutional networks for skeleton-based action recognition. In: \textit{AAAI}; 2018.
\bibitem{shi2019twoStream} Shi L, et al. Two-stream adaptive graph convolutional networks for skeleton-based action recognition. In: \textit{CVPR}; 2019.
\bibitem{zhu2023motionbert} Zhu W, et al. MotionBERT: a unified perspective on learning human motion representations. In: \textit{ICCV}; 2023.
\bibitem{mehraban2024motionagformer} Mehraban S, et al. MotionAGFormer: enhancing 3D human pose estimation with a transformer-GCNFormer network. In: \textit{WACV}; 2024.
\bibitem{winter2009biomechanics} Winter DA. \textit{Biomechanics and Motor Control of Human Movement}. 4th ed.\ Wiley; 2009.
\end{thebibliography}

\end{document}
